\documentclass[a4paper]{article}
\usepackage[margin=2cm]{geometry}
\usepackage{graphicx}
\usepackage{subcaption}
\usepackage{caption}
\usepackage{booktabs}
\usepackage{array}
\usepackage{tikz}
\usetikzlibrary{arrows.meta, positioning, calc}
\usepackage[backend=biber, style=numeric, sorting=none]{biblatex}
\addbibresource{refs.bib}

\captionsetup[figure]{skip=2pt}

\newcommand{\schemeline}[1]{%
  \noindent
  \begin{minipage}[t]{0.32\linewidth}
    \includegraphics[width=\linewidth]{residuals_#1.pdf}
  \end{minipage}\hfill
  \begin{minipage}[t]{0.32\linewidth}
    \includegraphics[width=\linewidth]{errors_cp_#1.pdf}
  \end{minipage}\hfill
  \begin{minipage}[t]{0.32\linewidth}
    \includegraphics[width=\linewidth]{errors_ch_#1.pdf}
  \end{minipage}
  \captionof{figure}{\texttt{#1}}
  \vspace{4pt}
}

\begin{document}

%% -----------------------------------------------------------------------
%% Page 1 -- Test case description
%% -----------------------------------------------------------------------
\begin{center}
  {\LARGE\bfseries Half-cylinder at Mach 17.605 -- Hybrid mesh study}
\end{center}

\bigskip

\noindent
\textbf{Test case.}
We study the hypersonic flow over a half-cylinder at Mach $17.605$, taken from the NASA
benchmark database~\cite{FUN3D} (see also~\cite{Delmas2026}, Chapter~6).  Reference solutions for the pressure coefficient~$C_p$
and the heat-flux coefficient~$C_h$ (Stanton number) are provided by the LAURA and FUN3D
solvers from NASA.  Adherence boundary conditions are applied on the cylinder wall with an
imposed wall temperature of $\theta_\mathrm{wall} = 500~\mathrm{K}$.
The free-stream conditions are listed in Table~\ref{tab:params}.

\bigskip

\begin{table}[h!]
  \centering
  \caption{Free-stream parameters for the half-cylinder at Mach $17.605$.}
  \label{tab:params}
  \renewcommand{\arraystretch}{1.3}
  \begin{tabular}{lll}
    \toprule
    Parameter & Symbol & Value \\
    \midrule
    Mach number           & $\mathrm{Ma}_\infty$    & $17.605$ \\
    Free-stream density   & $\rho_\infty$           & $10^{-3}~\mathrm{kg\,m^{-3}}$ \\
    Free-stream velocity  & $V_\infty$              & $5000~\mathrm{m\,s^{-1}}$ \\
    Free-stream pressure  & $p_\infty$              & $57.616~\mathrm{Pa}$ \\
    Free-stream temperature & $\theta_\infty$       & $200~\mathrm{K}$ \\
    Wall temperature      & $\theta_\mathrm{wall}$  & $500~\mathrm{K}$ \\
    Prandtl number        & $\mathrm{Pr}$           & $0.71$ \\
    Reynolds number       & $Re_L$                  & $3.7693\times10^5~\mathrm{m^{-1}}$ \\
    Dynamic viscosity     & $\mu$                   & Sutherland's law \\
    Heat conductivity     & $\kappa$                & From Prandtl number \\
    \bottomrule
  \end{tabular}
\end{table}

\bigskip

\noindent
\textbf{Mesh.}
All computations presented here use a \textbf{hybrid} mesh: the boundary layer is
discretised with quadrilateral (quad) cells aligned with the wall, while the rest of
the domain (far field) is filled with unstructured triangles.
The boundary layer zone contains 156 cells in the tangential direction and 24 cells in
the normal direction.
This hybrid configuration is more representative of practical unstructured meshes
encountered in industrial applications than a fully structured quad mesh.

\bigskip

\noindent
\textbf{Compared schemes.}
All schemes are subface-based, cell-centered finite volume approximate Riemann solvers
developed in~\cite{Delmas2025,Delmas2026}.
They share the same positivity-preserving wave-speed computation (ensuring positive
intermediate densities and internal energies) and differ only in how they treat the
contact wave and tangential velocity.

\begin{description}
  \item[\texttt{two\_wave}] (HLL-like)
    A single intermediate star state is assumed.
    Both contact discontinuities and shear layers are diffused, which leads to excessive
    numerical viscosity in the boundary layer and a significant overestimate of the wall
    heat flux at first order~\cite{Delmas2026}.

  \item[\texttt{three\_wave}] (HLLC-like)
    A central contact wave separates two star regions with distinct densities and
    tangential velocities.
    Shear layers and contact discontinuities aligned with the mesh are resolved exactly.
    However, this scheme is susceptible to carbuncle-type instabilities~\cite{Delmas2025},
    which appear strongly on unstructured and hybrid meshes.

  \item[\texttt{modified\_three\_wave}] (HLLCM-like)
    The three-wave density structure is retained but the tangential velocity is averaged
    as in the two-wave scheme.
    This suppresses carbuncle instabilities while preserving sharp contact resolution.
    However, the residual viscosity on shear layers does not vanish with normal-direction
    mesh refinement on anisotropic elements, preventing convergence to the correct heat
    flux~\cite{Delmas2026}.

  \item[\texttt{multi\_point\_iso}] (Isotropic multi-point)
    A node-conservative scheme that couples all Riemann problems sharing a mesh node
    through a local linear system, introducing genuine multidimensional coupling.
    The nodal velocity is computed with isotropic weights that cancel face-area
    anisotropy, so that the numerical diffusion on shear layers decreases when the mesh
    is refined in the wall-normal direction without growing when refined along the wall.
    This scheme is stable (no carbuncle) and accurately predicts heat fluxes on
    anisotropic hybrid meshes~\cite{Delmas2025,Delmas2026}.
\end{description}

\bigskip

\newpage

\noindent
\textbf{Parametric study.}
The sensitivity of the four Eulerian schemes (\texttt{two\_wave}, \texttt{three\_wave},
\texttt{modified\_three\_wave}, \texttt{multi\_point\_iso}) to the mesh resolution is
assessed by varying two independent parameters:
\begin{itemize}
  \item \textbf{First cell size} $d_\mathrm{wall}$ (x-axis): size of the first
        quadrilateral cell at the wall in the boundary layer, ranging from
        $2\times10^{-5}$ to $10^{-3}$ (10 values, logarithmically spaced).
  \item \textbf{Far-field cell size} (y-axis): characteristic size of the triangular
        cells in the far-field region, ranging from $0.2$ to $2.0$ (10 values,
        linearly spaced).
\end{itemize}
This produces a $10\times10$ grid of 100 simulations per scheme.
For each simulation, the RMSE relative error (\%) on $C_p$ and $C_h$ with respect to
the NASA reference solution and the normalised final residual $R_\mathrm{final}/R_0$
are collected and displayed on the maps on the following pages.

\bigskip

\noindent
\textbf{Reading the maps.}
The diagram on the following page illustrates how the mesh resolution varies across the
parameter space.
Each map covers the same $10\times10$ grid; a point near a given corner corresponds to
the mesh configuration described by the arrows.
For each corner, the left image shows the full mesh and the right image is a close-up
of the boundary layer.

\begin{center}
\resizebox{\textwidth}{!}{%
\begin{tikzpicture}[
    >=Stealth,
    every node/.style={font=\small},
    box/.style={draw, minimum width=4cm, minimum height=3cm, fill=white},
    img/.style={draw=gray!60, thick, inner sep=1pt}
  ]

  % Box: fond blanc, cellules de coin colorées, grille 10x10, bordure
  \fill[white] (-2,-1.5) rectangle (2,1.5);
  \fill[blue!15] (-2,1.2) rectangle (-1.6,1.5);
  \fill[blue!15] (1.6,1.2) rectangle (2,1.5);
  \fill[blue!15] (-2,-1.5) rectangle (-1.6,-1.2);
  \fill[blue!15] (1.6,-1.5) rectangle (2,-1.2);
  \draw[gray!50, very thin, xstep=0.4, ystep=0.3] (-2,-1.5) grid (2,1.5);
  \node[draw, minimum width=4cm, minimum height=3cm, inner sep=0, fill=none] (map) at (0,0) {};

  \node[anchor=north, font=\footnotesize\itshape] at (map.south)
    {First cell size $d_\mathrm{wall}$};
  \node[rotate=90, anchor=south, font=\footnotesize\itshape] at (map.west)
    {Far-field cell size};

  % Points aux centres des cellules de coin
  \filldraw[black] (-1.8, 1.35) circle (2pt);
  \filldraw[black] ( 1.8, 1.35) circle (2pt);
  \filldraw[black] (-1.8,-1.35) circle (2pt);
  \filldraw[black] ( 1.8,-1.35) circle (2pt);

  %% Top-left: coarse FF, fine BL  (above box, left side)
  \node[img] (tl) at (-3.5, 4.5) {%
    \begin{tabular}{@{}c@{\hspace{2pt}}c@{}}
      \includegraphics[width=3cm, trim={2434bp 1200bp 2436bp 1200bp}, clip]%
        {corners_visu/top_left_mesh.png} &
      \includegraphics[width=3cm, trim={3600bp 2442bp 3200bp 2442bp}, clip]%
        {corners_visu/top_left_mesh.png}
    \end{tabular}};
  \node[font=\tiny\bfseries, align=center, anchor=north] at (tl.south)
    {Coarse FF / Fine BL \quad $f_\infty=2.0$,\ $d_w=2{\times}10^{-5}$};
  \draw[->] (-1.8, 1.35) -- (tl.south);

  %% Top-right: coarse FF, coarse BL  (above box, right side)
  \node[img] (tr) at (3.5, 4.5) {%
    \begin{tabular}{@{}c@{\hspace{2pt}}c@{}}
      \includegraphics[width=3cm, trim={2434bp 1200bp 2436bp 1200bp}, clip]%
        {corners_visu/top_right_mesh.png} &
      \includegraphics[width=3cm, trim={3600bp 2442bp 3200bp 2442bp}, clip]%
        {corners_visu/top_right_mesh.png}
    \end{tabular}};
  \node[font=\tiny\bfseries, align=center, anchor=north] at (tr.south)
    {Coarse FF / Coarse BL \quad $f_\infty=2.0$,\ $d_w=10^{-3}$};
  \draw[->] ( 1.8, 1.35) -- (tr.south);

  %% Bottom-left: fine FF, fine BL  (below box, left side)
  \node[img] (bl) at (-3.5, -4.5) {%
    \begin{tabular}{@{}c@{\hspace{2pt}}c@{}}
      \includegraphics[width=3cm, trim={2434bp 1200bp 2436bp 1200bp}, clip]%
        {corners_visu/bottom_left_mesh.png} &
      \includegraphics[width=3cm, trim={3600bp 2442bp 3200bp 2442bp}, clip]%
        {corners_visu/bottom_left_mesh.png}
    \end{tabular}};
  \node[font=\tiny\bfseries, align=center, anchor=south] at (bl.north)
    {Fine FF / Fine BL \quad $f_\infty=0.2$,\ $d_w=2{\times}10^{-5}$};
  \draw[->] (-1.8,-1.35) -- (bl.north);

  %% Bottom-right: fine FF, coarse BL  (below box, right side)
  \node[img] (br) at (3.5, -4.5) {%
    \begin{tabular}{@{}c@{\hspace{2pt}}c@{}}
      \includegraphics[width=3cm, trim={2434bp 1200bp 2436bp 1200bp}, clip]%
        {corners_visu/bottom_right_mesh.png} &
      \includegraphics[width=3cm, trim={3600bp 2442bp 3200bp 2442bp}, clip]%
        {corners_visu/bottom_right_mesh.png}
    \end{tabular}};
  \node[font=\tiny\bfseries, align=center, anchor=south] at (br.north)
    {Fine FF / Coarse BL \quad $f_\infty=0.2$,\ $d_w=10^{-3}$};
  \draw[->] ( 1.8,-1.35) -- (br.north);

\end{tikzpicture}}
\end{center}

%% -----------------------------------------------------------------------
%% Following pages -- one row per scheme: residual | Cp error | Ch error
%% -----------------------------------------------------------------------

\newpage
\schemeline{two\_wave}
\schemeline{three\_wave}
\schemeline{modified\_three\_wave}
\schemeline{multi\_point\_iso}
\schemeline{ZB\_AMISO\_LVPPP}

\newpage
\printbibliography

\end{document}
